\begin{figure}[t]
\centering
\begin{tikzpicture}[>=Stealth, line width=0.9pt, font=\footnotesize]

% ---- axes ----
\draw[->,gray!60] (0,0) -- (13.2,0) node[above left=1pt and 0pt,black]{time};
\draw[->,gray!60] (0,0) -- (0,5.2) node[above,black,align=center]{adaptive\\capacity};

% ---- disruptions: live incidents (red) and simulated probes (audit / red-team, gray) ----
% D1 probe (x=3), D2 incident (x=5.2), D3 probe (x=7.4)
\foreach \x in {3,5.2,7.4} { \draw[gray!35,dotted] (\x,0.1) -- (\x,4.7); }
\node[gray!75,font=\scriptsize,align=center] at (3,-0.45)   {audit /\\red-team};
\node[red!70!black,font=\scriptsize,align=center] at (5.2,-0.45) {live\\incident};
\node[gray!75,font=\scriptsize,align=center] at (7.4,-0.45) {red-team};
\node[gray!85,font=\scriptsize] at (3,4.95) {\scriptsize(simulated disruption)};

% disruption markers on the axis
\node[gray!75] at (3,0.0)   {\scriptsize$\triangledown$};
\node[red!70!black] at (5.2,0.0) {\scriptsize$\blacktriangledown$};
\node[gray!75] at (7.4,0.0) {\scriptsize$\triangledown$};

% ---- trajectories: all dip at every disruption (real OR simulated); recovery differs ----
% green: monitoring + systematic learning -> each shock converted into a net gain
\draw[green!55!black]
  (1.8,3)--(3,3)--(3.2,2.3)--(4.0,3.5)--(5.2,3.5)--(5.4,2.8)--(6.2,4.0)--(7.4,4.0)--(7.6,3.3)--(8.4,4.6)--(9.6,4.75);
% blue: responds to every probe and incident, recovers to baseline, learns nothing lasting
\draw[blue!45!black]
  (1.8,3)--(3,3)--(3.2,2.3)--(4.0,3.0)--(5.2,3.0)--(5.4,2.3)--(6.2,3.0)--(7.4,3.0)--(7.6,2.3)--(8.4,3.0)--(9.6,3.0);
% red: no learning loop -> each shock leaves less capacity than before
\draw[red!75!black]
  (1.8,3)--(3,3)--(3.2,2.3)--(4.0,2.6)--(5.2,2.6)--(5.4,1.9)--(6.2,2.3)--(7.4,2.3)--(7.6,1.6)--(8.4,1.9)--(9.6,1.6);

% ---- mechanism annotations on the curves ----
\node[green!50!black,font=\scriptsize,align=center] at (6.0,4.75) {monitor + learn:\\each probe/incident\\becomes a gain};
\draw[green!50!black,->] (6.0,4.45) to[bend right=10] (5.7,3.95);
\node[red!65!black,font=\scriptsize,align=center] at (6.55,1.25) {no learning loop:\\lessons lost,\\capacity erodes};
\draw[red!65!black,->] (6.55,1.55) to[bend left=10] (6.1,2.25);

% ---- endpoint labels (outcome): all recover, but to different levels ----
\node[green!55!black,right,font=\scriptsize,align=left] at (9.6,4.75) {recovers\\\emph{stronger}};
\node[blue!45!black,right,font=\scriptsize,align=left]  at (9.6,3.0)  {recovers to\\baseline};
\node[red!75!black,right,font=\scriptsize,align=left]   at (9.6,1.6)  {recovers\\\emph{weaker}};

% ---- early ambiguity note ----
\node[align=center,font=\scriptsize] at (1.05,1.5)
  {one probe is\\ambiguous: all three\\dip and recover\\alike at first};
\draw[->,gray] (1.05,2.15) -- (2.55,2.92);

% ---- brace: longitudinal window ----
\draw[decorate,decoration={brace,amplitude=5pt,mirror},gray!70]
  (2.6,-1.05) -- (9.5,-1.05)
  node[midway,below=4pt,black,font=\scriptsize]
  {longitudinal tracing across disruptions: divergence becomes observable};

\end{tikzpicture}
\caption{\textbf{Monitoring and systematic learning are what separate the
trajectories.} Disruptions are both real (live incidents) and simulated
(audits, certifications, red-team exercises); all are probes that every
organization dips under and recovers from, so any single probe is ambiguous.
All three \emph{recover} after each disruption; what differs is the level they
recover to. Where monitoring feeds a systematic post-incident learning loop,
each shock is converted into a net gain and the organization recovers
\emph{stronger} (an antifragile path); where it responds but institutionalizes
nothing, it recovers only \emph{to baseline} (engineering-style robustness, no net
change); and where no learning loop exists, lessons are lost and it recovers
\emph{weaker}, growing progressively more brittle. Only tracing the loop across
successive disruptions makes the difference---and thus resilience versus luck,
capability versus compliance---observable.}
\label{fig:trace_vs_snapshot}
\end{figure}
